% TOP_PROBLEM: {"id": 7200014, "title": "Wikipedia geometry item 14: Nagata–Biran conjecture that if $X$ is a smooth algebraic surface and $L$ is an ample line bund…", "category_id": 5, "category": {"id": 5, "name": "algebraic_geometry", "display_name": "Algebraic Geometry", "description": "Geometric objects defined by polynomial equations.", "slug": "algebraic-geometry", "order_index": 5, "created_at": "2026-07-31T15:26:25.671Z"}, "status": "open"}
% FIRST_POSTED: null
% ENTRY_KIND: statement_only
% !TeX program = lualatex
\documentclass[11pt]{article}
\usepackage[margin=25mm]{geometry}
\usepackage{fontspec}
\setmainfont{Latin Modern Roman}
\usepackage{amsmath,amssymb,mathtools}
\usepackage{unicode-math}
\setmathfont{Latin Modern Math}
\usepackage{xurl,hyperref}
\hypersetup{colorlinks=true,urlcolor=blue}
\setcounter{secnumdepth}{0}
\setlength{\parindent}{0pt}
\setlength{\parskip}{5pt}
\setlength{\emergencystretch}{3em}
\newcommand{\R}{\mathbb R}
\newcommand{\N}{\mathbb N}
\newcommand{\Z}{\mathbb Z}
\newcommand{\Q}{\mathbb Q}
\newcommand{\C}{\mathbb C}
\newcommand{\D}{\mathbb D}
\newcommand{\F}{\mathbb F}
\newcommand{\sref}[2]{\hyperlink{source:#1:#2}{[S#2]}}
\newcommand{\cataloguescope}[1]{#1}
\begin{document}
% UnsolvedMath ID 7200014; source catalogue code AMR-071-0014
\setcounter{section}{7200013}
\section{\texorpdfstring{Wikipedia geometry item 14: Nagata–Biran conjecture that if $X$ is a smooth algebraic surface and $L$ is an ample line bund…}{Wikipedia geometry item 14: Nagata–Biran conjecture that if X is a smooth algebraic surface and L is an ample line bund…}}\label{7200014}
{\small UnsolvedMath ID 7200014 · Algebraic Geometry}
\subsection{Definitions and mathematical statement}

\paragraph{Shared notation.}$\mathbb N$, $\mathbb Z$, $\mathbb Q$, $\mathbb R$, and $\mathbb C$ denote the natural numbers, integers, rationals, reals, and complex numbers, respectively. All additional parameters and hypotheses are specified in the question.


\paragraph{Statement recorded in the supplied source.}
Nagata–Biran conjecture that if $X$ is a smooth algebraic surface and $L$ is an ample line bundle on $X$ of degree $d$, then for sufficiently large $r$, the Seshadri constant satisfies $\varepsilon(p_1,\ldots,p_r;X,L) = d/\sqrt{r}$.
\par\smallskip\textit{Formulation references:} \sref{7200014}{1}, \sref{7200014}{2}.
\subsection{Short English statement}
Nagata–Biran conjecture that if $X$ is a smooth algebraic surface and $L$ is an ample line bundle on $X$ of degree $d$, then for sufficiently large $r$, the Seshadri constant satisfies $\varepsilon(p_1,\ldots,p_r;X,L) = d/\sqrt{r}$.
\subsection{Sources}
\begin{itemize}
\item[\hypertarget{source:7200014:1}{S1}] ULAM.AI and the UnsolvedMath Contributors, UnsolvedMath: A Curated Collection of Open Mathematics Problems, record 7200014 (AMR-071-0014). Catalogue attribution under CC BY 4.0.
\url{https://huggingface.co/datasets/ulamai/UnsolvedMath}.
\item[\hypertarget{source:7200014:2}{S2}] Original problem formulation and mathematical context.
\url{https://en.wikipedia.org/w/index.php?title=List_of_unsolved_problems_in_mathematics&oldid=1366636767#Geometry}.
\end{itemize}
\label{end:7200014}
\end{document}
